% Section 6.1, from lectures/L06/appendix/a_watchdog.md.

\section{The Watchdog}
\label{c6:sec:watchdog}

\subsection{What it actually detects}
\label{c6:sec:detects}
A watchdog is a timer that resets the device unless your program keeps telling it not to. That is
the whole mechanism, and everything useful and everything disappointing about it follows from
that one sentence.

\textbf{It detects exactly one thing: that your program stopped reaching the instruction that pets
it.}

So it catches a program stuck in a loop that never exits, a program blocked waiting for an
interrupt that will never arrive, and a program that jumped somewhere meaningless and is executing
whatever it found. Those are real failures and they are the ones hardest to catch any other way.

Here is what it does \textbf{not} catch, and the list is longer:
\begin{itemize}
  \item \textbf{A program computing wrong answers quickly.} It is petting the watchdog on schedule
    and is completely broken.
  \item \textbf{The atomicity bug from Chapter~3.} The main loop reads a torn 16-bit value and
    carries on. Nothing is stuck.
  \item \textbf{A stack that has grown into your variables}, right up until it corrupts a return
    address and the program jumps somewhere that stops petting (\secref{c4:sec:overflow}).
  \item \textbf{A driver writing to the wrong port register.} The program runs perfectly and the
    hardware does the wrong thing.
  \item \textbf{Anything about correctness at all.}
\end{itemize}

And it introduces a failure of its own: \textbf{a program that pets the watchdog from inside the
wrong place.} A handler that runs on a timer and pets the watchdog will keep petting it while the
main loop is completely stuck, which converts a watchdog from a safety mechanism into decoration.
Pet it from one place, in the main loop, at the point that proves the whole loop ran.

\subsection{The two modes}
\label{c6:sec:modes}
The watchdog can do one of two things when it times out, and it is worth being clear that these
are different tools.

\textbf{System reset}, with \code{WDE} set. The device restarts as though it had been power
cycled, except that SRAM keeps its contents and \code{MCUSR} records that the watchdog was the
cause. This is the one you want in a product.

\textbf{Interrupt}, with \code{WDIE} set. The timeout raises the \code{WDT} interrupt instead, and
your handler runs. Nothing is reset. This is useful for waking from sleep (\secref{c6:sec:sleep})
and for recording diagnostics before a deliberate reset, and it is \textbf{not} a safety mechanism,
because a program too broken to pet the watchdog is usually too broken to run a handler sensibly.

Both bits can be set at once, which the datasheet calls interrupt-and-reset mode: the first
timeout interrupts, and because the hardware clears \code{WDIE} on entry, the second one resets.
That gives you one chance to save something before the restart.

\subsection{The registers}
\label{c6:sec:registers}

\bookfigure{watchdog_registers}{\code{WDTCSR} drawn as its eight bits, with \code{WDP3} at bit 5
and \code{WDP2} to \code{WDP0} at bits 2 to 0, and \code{MCUSR} below it with \code{WDRF} marked
as the bit set when the watchdog caused the last reset.}{c6:fig:registers}

\begin{booktable}{@{}p{5.2em}p{4.4em}L@{}}
  \thead{Register} & \thead{Address} & \thead{Holds} \\
  \midrule
  \code{WDTCSR} & \code{0x60} & The timeout selector, \code{WDE}, \code{WDIE}, and the window bit
    \code{WDCE}. \\
  \code{MCUSR} & \code{0x54} & Why the last reset happened: power-on, external, brown-out, or
    watchdog. \\
\end{booktable}

\textbf{\code{WDTCSR} is in extended I/O and \code{MCUSR} is not}, so the first needs \code{lds} and
\code{sts} and the second can use \code{in} and \code{out}. Two registers in one subroutine, two
addressings.

\textbf{The timeout selector is four bits that are not next to each other.} \code{WDP3} is bit 5,
and \code{WDP2} to \code{WDP0} are bits 2 to 0, with \code{WDCE} and \code{WDE} in between. The
datasheet numbers the timeouts 0 to 9, and \textbf{that number is not the byte you write}:

\begin{narrowtable}{@{}lcc@{}}
  \thead{Timeout} & \thead{Selector} & \thead{Byte} \\
  \midrule
  16\,ms & 0 & \code{0x00} \\
  125\,ms & 3 & \code{0x03} \\
  2\,s & 7 & \code{0x07} \\
  4\,s & 8 & \textbf{\code{0x20}} \\
  8\,s & 9 & \textbf{\code{0x21}} \\
\end{narrowtable}

\textbf{Writing \code{0x08} because you wanted four seconds is the most expensive typo in this
chapter.} \code{0x08} is \code{WDE} with the selector left at zero, so it means ``reset me if I go
quiet for sixteen milliseconds''. For most programs that is immediately, and for ever: the device
resets, runs a few hundred microseconds of startup, and resets again. Reprogramming it then means
winning a race against the watchdog, which is possible and unpleasant.

Nobody remembers that table, and nobody should try to. What is worth remembering is its shape:
the four selector bits are not contiguous, so the byte is never the timeout number.

\subsection{The timed write sequence}
\label{c6:sec:timed}
The watchdog is the one peripheral on this device that refuses to be configured carelessly, and
the reason is that a program which has gone wrong must not be able to switch off the thing that
would rescue it.

\bookfigure[htbp]{timed_write}{Six steps in order: save \code{SREG} and \code{cli}, \code{wdr}, clear
\code{WDRF} in \code{MCUSR}, then two shaded writes that must fall within four cycles of each
other, then restoring \code{SREG}.}{c6:fig:timed}

So changing \code{WDE} or the timeout takes two writes:
\begin{enumerate}
  \item Write \code{WDCE} and \code{WDE} \textbf{together}. This opens a window.
  \item Within \textbf{four cycles}, write the value you actually want, with \code{WDCE} clear.
\end{enumerate}

Four cycles, not four instructions. Two consecutive \code{sts} instructions are two cycles each,
so they fit with one cycle to spare: a single \code{nop} between them still works, and anything
longer does not.

\textbf{This is why the sequence starts with \code{cli}.} An interrupt landing between the two
writes takes at least six cycles to reach your handler and then runs the whole thing
(\secref{c3:sec:sequence}). The window would be long gone. Turning interrupts off is not defensive
style here; it is the only way the sequence can be made to work at all.

\textbf{And it is why \code{WDRF} must be cleared first.} The hardware will not let you clear
\code{WDE} while \code{WDRF} is set in \code{MCUSR}. That rule exists so that a device which has
just been reset by the watchdog cannot immediately disable it and carry on being broken. The
consequence for you is specific: a driver that forgets to clear \code{WDRF} works perfectly the
first time it runs and cannot switch the watchdog off after a watchdog reset, which is precisely
when you most want to.

\subsubsection{What a failed sequence leaves behind}
If the second write misses its window, the hardware ignores it. \code{WDCE} clears itself after
four cycles, and what remains is whatever the \emph{first} write was allowed to change on its own:
\code{WDE} is set, because setting it never needs the window, and the timeout selector is where it
was, because changing it does. On a device fresh from reset the selector is zero, so the register
holds \code{0x08}.

That is system reset mode with the shortest timeout. \textbf{A sequence that is merely slightly
too slow does not fail safe; it fails into system reset mode, and from reset into the most
aggressive setting the register has.} The simulator models this exactly, and the test suite checks
for it: if \code{WatchdogDriver.InitWritesTheByteAsked} ever reports \code{WDTCSR} holding
\code{0x08}, that is what happened.

\subsection{Choosing a timeout}
\label{c6:sec:timeout}
Ten timeouts, from 16\,ms to 8\,s, roughly doubling. The choice is a genuine engineering decision
with a cost in both directions.

\textbf{Too short} and the watchdog fires during normal operation, on the one iteration that took
longer than usual. A program that resets itself occasionally under load is far worse than one that
does not have a watchdog at all, because now you are debugging the watchdog as well.

\textbf{Too long} and a program that has stopped stays stopped for that long. Eight seconds is a
long time for a machine to be doing nothing.

The rule is: \textbf{the shortest timeout that comfortably exceeds your worst-case loop time}.
Worst case, not typical: the iteration that reads a sensor, and handles two interrupts, and takes
the slow branch.

\textbf{And the timeouts are nominal.} They come from an internal oscillator that runs at roughly
128\,kHz and varies with supply voltage and temperature. The datasheet's ``1.0\,s'' and the
``1024\,ms'' you will find in other people's code are the same setting described two ways, and
neither is a promise. Leave real margin, and do not build anything on the exact value.
